
\begin{abstract}
本文针对微构体中导电介质填充的仿真优化问题，提出了基于三维几何建模、蒙特卡洛模拟和并查集连通性判定的系统解决方案。

\textbf{针对问题一：}本问是单颗粒导通判定问题，需判断给定坐标的单个金属A圆柱体能否使微构体导通。首先，建立了圆柱体的三维几何模型，推导了考虑周期性边界条件的圆柱体表面到电极面的最短距离公式。接着，对附件1中12个颗粒逐一计算到两侧电极的最短表面距离。结果讨论显示，所有12个配置均无法导通，单个长度5000nm的圆柱体在边长10000nm的微构体中无法同时触及两侧电极面，\textbf{导通概率为0\%}。

\textbf{针对问题二：}本问是多颗粒导通概率计算问题，需计算不同填充率下的导通概率。首先，建立了三维线段间最短距离的几何算法和基于并查集的连通性判定框架。接着，引入空间哈希加速近邻搜索，大幅降低计算复杂度。然后，对填充率0.50\%、0.60\%、0.70\%和1.00\%分别运行200次蒙特卡洛模拟。结果讨论显示，\textbf{四种填充率下的导通概率均为100\%}，这是由于金属A圆柱体的长径比高达83:1，渗透阈值远低于0.50\%。对附件2中49个颗粒的固定配置进行分析，填充率仅为0.0693\%，\textbf{该配置未导通}。

\textbf{针对问题三：}本问是最小填充率确定问题，需找到使导通概率不低于90\%的最小金属A填充率。首先，分析问题二的渗透现象确定搜索区间为5至80个颗粒。接着，采用二分搜索策略，每步运行200次蒙特卡洛模拟估计导通概率。经过6轮迭代收敛至最优解，并进行500次大样本验证。结果讨论显示，使导通概率不低于90\%的\textbf{最小填充率为0.0551\%}，对应\textbf{39个金属A圆柱体}，验证导通概率为90.60\%。

\textbf{针对问题四：}本问是双金属成本优化问题，需在导通概率不低于90\%的约束下最小化总成本。首先，计算金属A和金属B的单颗粒成本分别为0.01484元和0.00168元。接着，测试纯金属B方案发现需约2000个球体才能达到99\%导通概率，成本高达3.35元。进一步分析表明，金属B球体的渗透效率远低于金属A圆柱体，任何以金属B替代金属A的混合方案均因球体的低效而导致总成本上升。\textbf{最优方案为纯金属A}，数量39个，\textbf{最低成本为0.579元}。

综上，本文为微构体导电介质填充问题提供了从几何建模到统计模拟再到优化决策的完整解决方案，揭示了高长径比圆柱形填料具有超低渗透阈值的物理特性。

\keywords{微构体；导电介质；蒙特卡洛模拟；渗透阈值；并查集；三维几何建模}
\label{abstract:end}
\end{abstract}
